\section{Problem and relevance}
\label{sec:intro}
Dengue is endemic in Singapore, with substantial year-to-year variation and recurring
epidemics. An early warning could help public-health teams prioritize vector control,
risk communication and service readiness. Singapore forecasting studies have used
climate and surveillance information at longer horizons
\citep{hii2012forecast,shi2016three}; here we ask a narrower predictive question:
\emph{for Singapore dengue surveillance teams, can case counts and weather available by
week $t$ identify a materially elevated DF week at $t+h$?}

We compare three operationally distinct outcomes: crossing a recent-count threshold,
crossing a historical seasonal threshold, and a relative increase with a minimum count.
The onset variants ask the more decision-relevant question of whether an elevated period
can be flagged before the series is already above its threshold. We then assess whether
weather improves on recent cases and transparent seasonal or persistence baselines.
The analysis is predictive, not causal: an association between weather and future counts
does not identify the effect of weather or the impact of an intervention.
